\scanpage{97}
% manuscript lines 1–2
\textcircled{2} Supposons qu'on ait $=$
\[
\text{b'')}\qquad v^n\ \text{est bijectif}
\]

% manuscript lines 3–6
Alors $\psi$ est \uline{injectif}. \deleted{Alors} De plus, les $B^i$ sont plats [resp. \deleted{libres,} projectifs] sur $K$ sss les $A^i,C^i$ le sont.

% manuscript lines 2–11
% slanted
\margin{\uline{N.B.} On devrait vérifier que (b'') ne dépend pas du choix des $u^n$. [Si les $A^i,B^i$ de type fini, celà se ramène au cas d'un corps de base, (comme dans $3^\circ$) grâce à Nakayama]}

% manuscript lines 7–10
\textcircled{3} \deleted{Réciproquement, supposons \deleted{\unsure{les|que}} les $A^i,B^i,C^i$ projectifs de type fini [= plats de présentation finie]} Pour tt $x\in\operatorname{Spec} K$, soit

% manuscript lines 11
\[
 P_{A,x}\deleted{\unsure{A}}(t)=\sum\operatorname{rg}_{k(x)}(A^i(x))t^i\in\mathbb Z[[t]]
\]

% manuscript lines 12–13
et définissons de $=$ $P_{B,x}(t)$, $P_{C,x}(t)$. Alors on a \deleted{\unsure{tt}} pour tt $x$ [moyennant a]

% manuscript lines 14
\[
P_{B,x}\ll P_{A,x}\,P_{C,x}\tag*{\textcircled{*}}
\]

% manuscript lines 15–18
[majoration coefficient par coefficient]. Si (b'') est vérifié, alors on a $=$

% manuscript lines 16–17
\margin{c) $P_{B,x}=P_{A,x}P_{C,x}$}

% manuscript lines 17–19
Si les $A^i,B^i,C^i$ sont plats de prés. finie [= projectifs de type fini], alors c) implique aussi b''), \deleted{donc \unsure{Donc}}
% archive sheet number: 96

